\documentclass{zkdl-presentation-template}

\title[Linear Algebra]{\textbf{Mathematics for Cryptography III: Linear Algebra}}
\author{Blockstream Research}
\date{\today}
\homepage{research.blockstream.com}
\github{BlockstreamResearch/zk-book}

\begin{document}
    \maketitleframe

    \begin{frame}{Plan}
      \tableofcontents[hideallsubsections]
    \end{frame}


    \section{Motivation}

    \begin{frame}{Why Linear Algebra?}
        In cryptography we commonly need to work with a large set of \textit{linear} relations, e.g. $\alpha_1x_1 + \alpha_2x_2 + \dots + \alpha_nx_n = 0$ over a field $\mathbb{F}$.
        \pause

        We collect variables into \textbf{vectors} (in bold): $\mathbf{x}=(x_1,\dots,x_n)$, and relations into \textbf{matrices}:
        \begin{equation*}
            \begin{cases}
                \alpha x_1 + \beta x_2 = 0 \\
                \gamma x_1 + \delta x_3 = 0 \\
                \varepsilon x_2 + \zeta x_3 = 0
            \end{cases}
            \quad \Longleftrightarrow \quad
            A\mathbf{x} = \mathbf{0}, \quad A = \begin{bmatrix}
                \alpha & \beta & 0 \\
                \gamma & 0 & \delta \\
                0 & \varepsilon & \zeta
            \end{bmatrix}.
        \end{equation*}
        \pause

        \begin{remarkblock}{Where we use it}
            Linear PCPs and R1CS (Hadamard products of large vectors), inner-product arguments (Bulletproofs), and lattices (short solutions of $A\mathbf{x} = \mathbf{0}$).
        \end{remarkblock}
    \end{frame}

    \begin{frame}{Geometric View}
        \begin{columns}[c]
            \begin{column}{0.33\textwidth}
                \centering
                \includegraphics[width=\textwidth]{../contents/images/linear-algebra/3d-system-1.png}\\
                \small\textcolor{oc-indigo-7}{\textit{(a)}} vectors span $\mathbb{R}^3$: a solution always exists.
            \end{column}
            \pause
            \begin{column}{0.33\textwidth}
                \centering
                \includegraphics[width=\textwidth]{../contents/images/linear-algebra/3d-system-2.png}\\
                \small\textcolor{oc-indigo-7}{\textit{(b)}} vectors on a plane: infinitely many or no solutions.
            \end{column}
            \pause
            \begin{column}{0.33\textwidth}
                \centering
                \includegraphics[width=\textwidth]{../contents/images/linear-algebra/3d-system-3.png}\\
                \small\textcolor{oc-indigo-7}{\textit{(c)}} vectors on a line: a solution exists only on this line.
            \end{column}
        \end{columns}

        \vspace{8pt}
        \centering
        Solutions of $\mathbf{v}_1x_1 + \mathbf{v}_2x_2 + \mathbf{v}_3x_3 = \mathbf{t}$ depend on how $\mathbf{v}_1,\mathbf{v}_2,\mathbf{v}_3$ are related.
    \end{frame}

    \begin{frame}{Determinant: the Volume Spanned by Columns}
        For $A = (\mathbf{v}_1,\mathbf{v}_2,\mathbf{v}_3)$, $\det A$ is the (signed) volume of the parallelepiped spanned by the columns.
        \begin{figure}
            \centering
            \resizebox{\textwidth}{!}{%
            \begin{tikzpicture}[scale=0.75, >=Stealth]
\begin{axis}[
    at={(0,0)},
    width=5.2cm,
    height=5.2cm,
    view={120}{25},
    axis lines=center,
    axis line style={thick, oc-gray-5},
    xmin=0, xmax=4.2,
    ymin=0, ymax=3.8,
    zmin=0, zmax=3.5,
    ticks=none,
    xlabel style={
        at={(axis description cs:1.02,0.03)},
        anchor=west
    },
    ylabel style={
        at={(axis description cs:0.02,1.03)},
        anchor=south
    },
    zlabel style={
        at={(axis description cs:0.05,0.98)},
        anchor=south
    },
]
\addplot3[
    fill=oc-blue-1,
    draw=oc-blue-5,
    opacity=0.75
] coordinates {
    (0,0,0) (3,0,0) (3,2,0) (0,2,0) (0,0,0)
};

\addplot3[
    fill=oc-blue-1,
    draw=oc-blue-5,
    opacity=0.75
] coordinates {
    (0,0,2) (3,0,2) (3,2,2) (0,2,2) (0,0,2)
};

\addplot3[oc-blue-5, opacity=0.75] coordinates {(0,0,0) (0,0,2)};
\addplot3[oc-blue-5, opacity=0.75] coordinates {(3,0,0) (3,0,2)};
\addplot3[oc-blue-5, opacity=0.75] coordinates {(3,2,0) (3,2,2)};
\addplot3[oc-blue-5, opacity=0.75] coordinates {(0,2,0) (0,2,2)};
\addplot3[
    ->,
    ultra thick,
    oc-red-7
] coordinates {(0,0,0) (3,0,0)};

\addplot3[
    ->,
    ultra thick,
    oc-blue-7
] coordinates {(0,0,0) (0,2,0)};

\addplot3[
    ->,
    ultra thick,
    oc-green-7
] coordinates {(0,0,0) (0,0,2)};

\node[
    oc-red-7,
    font=\small
] at (axis cs:1.5,-0.35,0.2) {$\mathbf{v}_1$};

\node[
    oc-blue-7,
    font=\small
] at (axis cs:-0.7,1,0.2) {$\mathbf{v}_2$};

\node[
    oc-green-7,
    font=\small
] at (axis cs:-0.25,0.3,1.15) {$\mathbf{v}_3$};

\node[
    font=\large\bfseries,
    anchor=south
] at (axis description cs:0.5,1.08)
{$\det(A)\neq0$};

\node[
    font=\small,
    anchor=north
] at (axis description cs:0.5,-0.15)
{3-dimensional span};

\end{axis}
\begin{axis}[
    at={(6cm,0)},
    width=5.2cm,
    height=5.2cm,
    view={120}{25},
    axis lines=center,
    axis line style={thick, oc-gray-5},
    xmin=0, xmax=4.2,
    ymin=0, ymax=3.8,
    zmin=0, zmax=3.5,
    ticks=none,
    xlabel style={
        at={(axis description cs:1.02,0.03)},
        anchor=west
    },
    ylabel style={
        at={(axis description cs:0.02,1.03)},
        anchor=south
    },
    zlabel style={
        at={(axis description cs:0.05,0.98)},
        anchor=south
    },
]

\addplot3[
    surf,
    domain=0:5,
    y domain=0:5,
    samples=8,
    samples y=8,
    shader=flat,
    draw=oc-gray-5,
    fill=oc-gray-1,
    opacity=0.45,
    line join=round
]
({x},{y},{0});
\addplot3[
    ->,
    ultra thick,
    oc-red-7
] coordinates {(0,0,0) (2,0,0)};

\addplot3[
    ->,
    ultra thick,
    oc-blue-7
] coordinates {(0,0,0) (0,2,0)};

\addplot3[
    ->,
    ultra thick,
    oc-green-7
] coordinates {(0,0,0) (3,2,0)};

\node[
    oc-red-7,
    font=\small
] at (axis cs:1.5,-0.15,0.5) {$\mathbf{v}_1$};

\node[
    oc-blue-7,
    font=\small
] at (axis cs:-0.3,2,0.35) {$\mathbf{v}_2$};

\node[
    oc-green-7,
    font=\small
] at (axis cs:2,2,0) {$\mathbf{v}_3$};

\node[
    font=\large\bfseries,
    anchor=south
] at (axis description cs:0.5,1.08)
{$\det(A)=0$};

\node[
    font=\small,
    anchor=north
] at (axis description cs:0.5,-0.15)
{2-dimensional span (plane)};

\end{axis}
\begin{axis}[
    at={(12cm,0)},
    width=5.2cm,
    height=5.2cm,
    view={120}{25},
    axis lines=center,
    axis line style={thick, oc-gray-5},
    xmin=0, xmax=5.2,
    ymin=0, ymax=3.8,
    zmin=0, zmax=4.2,
    ticks=none,
    xlabel style={
        at={(axis description cs:1.02,0.03)},
        anchor=west
    },
    ylabel style={
        at={(axis description cs:0.02,1.03)},
        anchor=south
    },
    zlabel style={
        at={(axis description cs:0.05,0.98)},
        anchor=south
    },
]

\addplot3[
    very thick,
    dashed,
    gray
] coordinates {
    (0,0,0)
    (3,3,0)
};

\addplot3[
    very thick,
    dashed,
    gray
] coordinates {
    (3,3,3)
    (3,3,0)
};
\addplot3[
    ->,
    ultra thick,
    oc-green-7
] coordinates {(0,0,0) (4,4,4)};

\addplot3[
    ->,
    ultra thick,
    oc-blue-7
] coordinates {(0,0,0) (2.75,2.75,2.75)};

\addplot3[
    ->,
    ultra thick,
    oc-red-7
] coordinates {(0,0,0) (1.5,1.5,1.5)};

\node[
    oc-green-7,
    font=\small
] at (axis cs:4,3.75,4.4) {$\mathbf{v}_3$};

\node[
    oc-blue-7,
    font=\small
] at (axis cs:2.75,2.5,3.25) {$\mathbf{v}_2$};

\node[
    oc-red-7,
    font=\small
] at (axis cs:1.5,1.25,2) {$\mathbf{v}_1$};

\end{axis}

\end{tikzpicture}%
            }
        \end{figure}
        \pause
        \begin{remarkblock}{Takeaway}
            A plane or a line has no volume ($\det A = 0$): the equation $A\mathbf{x} = \mathbf{t}$ has either no or infinitely many solutions. If $\det A \neq 0$, there is always exactly one.
        \end{remarkblock}
    \end{frame}

    \section{Vector Spaces}

    \begin{frame}{Vector Space}
        \begin{definition}
            A \textbf{vector space} $V$ over the field $\mathbb{F}$ is an abelian group for addition $+$ together with a scalar multiplication $\mathbb{F} \times V \to V$, $(\lambda,\mathbf{x}) \mapsto \lambda \mathbf{x}$, such that for any $\mathbf{v},\mathbf{u} \in V$ and $\lambda,\mu \in \mathbb{F}$:
            \begin{columns}
                \begin{column}{0.45\textwidth}
                    \begin{itemize}
                        \item $\lambda(\mathbf{u}+\mathbf{v}) = \lambda \mathbf{u} + \lambda \mathbf{v}$
                        \item $(\lambda + \mu)\mathbf{v} = \lambda \mathbf{v} + \mu \mathbf{v}$
                    \end{itemize}
                \end{column}
                \begin{column}{0.45\textwidth}
                    \begin{itemize}
                        \item $(\lambda \mu)\mathbf{v} = \lambda(\mu \mathbf{v})$
                        \item $1\cdot \mathbf{v} = \mathbf{v}$
                    \end{itemize}
                \end{column}
            \end{columns}
        \end{definition}
        \pause

        \begin{example}
            $\mathbb{F}^n$ with coordinate-wise operations is a vector space. So is $V_3 := \{\mathbf{x} \in \mathbb{F}^n: x_1+x_2+\dots+x_n = 0\}$: it is closed under addition and scalar multiplication.
        \end{example}
        \pause

        \begin{remark}
            Polynomials $\mathbb{R}[T]^{\leq n}$ also form a vector space over $\mathbb{R}$.
        \end{remark}
    \end{frame}

    \begin{frame}{Linear Subspace and Span}
        \begin{definition}
            Let $V$ be a vector space over $\mathbb{F}$ and $W \subseteq V$. $W$ is a \textbf{linear subspace} of $V$ if for any $\mathbf{w}_1,\mathbf{w}_2 \in W$ and $\alpha,\beta \in \mathbb{F}$ one has $\alpha\mathbf{w}_1 + \beta\mathbf{w}_2 \in W$.
        \end{definition}
        \pause

        \begin{example}
            For vectors $\mathbf{v}_1,\dots,\mathbf{v}_n \in V$, the \textbf{span}
            \begin{equation*}
                \text{span}(\mathbf{v}_1,\dots,\mathbf{v}_n) := \{\lambda_1\mathbf{v}_1 + \dots + \lambda_n\mathbf{v}_n: \lambda_1,\dots,\lambda_n \in \mathbb{F}\}
            \end{equation*}
            is a linear subspace of $V$.
        \end{example}
        \pause

        \begin{questionblock}{Question}
            Show that the set of solutions $\mathbf{x} \in \mathbb{F}^n$ of $\mathbf{u}_1x_1 + \dots + \mathbf{u}_nx_n = \mathbf{0}$ is a linear subspace of $\mathbb{F}^n$.
        \end{questionblock}
    \end{frame}

    \begin{frame}{Subspaces: Examples}
        \begin{example}
            Let $n \geq 2$ and
            \begin{equation*}
                U := \{\mathbf{x} \in \mathbb{F}^n: x_1 = x_2 = 0\}, \quad W := \{\mathbf{x} \in \mathbb{F}^n: x_1 + x_2 = 0\}.
            \end{equation*}
            Both $U$ and $W$ are linear subspaces of $\mathbb{F}^n$, and $U$ is a subspace of $W$.
        \end{example}
        \pause

        \begin{questionblock}{Question}
            Let $V := \{\mathbf{z} \in \mathbb{F}^n: \mathbf{z} \odot (\mathbf{1}_n - \mathbf{z}) = \mathbf{0}\}$ (bit-vectors, see the Hadamard product below). When is $V$ a linear subspace of $\mathbb{F}^n$?
        \end{questionblock}
    \end{frame}

    \section{Basis and Dimension}

    \begin{frame}{Linear Independence}
        \begin{definition}
            Vectors $\mathbf{v}_1, \dots, \mathbf{v}_n \in V$ are \textbf{linearly independent} if the only scalars $\lambda_1, \dots, \lambda_n \in \mathbb{F}$ with
            \begin{equation*}
                \lambda_1 \mathbf{v}_1 + \lambda_2 \mathbf{v}_2 + \cdots + \lambda_n \mathbf{v}_n = \mathbf{0}
            \end{equation*}
            are $\lambda_1 = \cdots = \lambda_n = 0$. Otherwise, they are \textbf{linearly dependent}.
        \end{definition}
        \pause

        \begin{example}
            $\mathbf{v}_1 = (1, 0, 0)$, $\mathbf{v}_2 = (0, 1, 0)$, $\mathbf{v}_3 = (0, 0, 1)$ are linearly independent in $\mathbb{R}^3$.
        \end{example}
        \pause

        \begin{example}
            $\mathbf{u}_1 = (1, 2, 3)$ and $\mathbf{u}_2 = (2, 4, 6)$ are linearly dependent: $2\mathbf{u}_1 - \mathbf{u}_2 = \mathbf{0}$.
        \end{example}
    \end{frame}

    \begin{frame}{Basis and Dimension}
        \begin{definition}
            Vectors $\mathbf{e}_1,\dots,\mathbf{e}_n$ are a \textbf{basis} of $V$ if $\text{span}(\mathbf{e}_1,\dots,\mathbf{e}_n) = V$ and $\mathbf{e}_1,\dots,\mathbf{e}_n$ are linearly independent.
        \end{definition}
        \pause

        \begin{example}
            $\mathbf{e}_1=(1,0)$, $\mathbf{e}_2=(0,1)$ is the \textit{canonical} basis of $\mathbb{F}^2$, but $\mathbf{u}_1=(1,2)$, $\mathbf{u}_2=(2,1)$ is a basis of $\mathbb{R}^2$ as well. For $\mathbb{F}[T]^{\leq n}$, the canonical basis is $1, T, \dots, T^n$.
        \end{example}
        \pause

        \begin{lemma}
            Any two bases of a vector space have the same length.
        \end{lemma}
        \pause

        \begin{definition}
            The \textbf{dimension} $\dim V$ is the length of any basis of $V$. For example, $\dim \mathbb{F}^n = \dim \mathbb{F}[T]^{\leq n-1} = n$.
        \end{definition}
    \end{frame}

    \begin{frame}{Dimension: Questions}
        \begin{questionblock}{Question \#1}
            Let $V := \{(x,y,z) \in \mathbb{R}^3: x+y+z=0\}$. Compute $\dim V$.
        \end{questionblock}
        \pause

        \begin{questionblock}{Question \#2}
            Is $3,\; T+2,\; T^2+4$ a basis of $\mathbb{F}_5[T]^{\leq 2}$?
        \end{questionblock}
        \pause

        \begin{questionblock}{Question \#3}
            How many distinct $1$-dimensional linear subspaces (lines through the origin) exist in $\mathbb{F}_3^3$?
        \end{questionblock}
    \end{frame}

    \section{Matrices}

    \begin{frame}{Matrix Operations}
        A matrix $A \in \mathbb{F}^{m \times n}$ has $m$ rows and $n$ columns; we write $A = (a_{i,j})$.
        \begin{definition}[Matrix operations]
            \begin{itemize}
                \item \textbf{Addition} and \textbf{scalar multiplication} are element-wise.
                \item For $A \in \mathbb{F}^{m \times n}$ and $B \in \mathbb{F}^{n \times p}$, the \textbf{product} $C := AB \in \mathbb{F}^{m \times p}$ has elements $c_{i,j} := \sum_{\ell=1}^{n} a_{i,\ell}b_{\ell,j}$.
                \item The \textbf{transpose} $A^{\top} \in \mathbb{F}^{n \times m}$ has elements $A^{\top}_{i,j} = A_{j,i}$.
            \end{itemize}
        \end{definition}
        \pause

        \begin{example}
            \begin{equation*}
                A = \begin{bmatrix}
                    1 & 1 & 2 \\
                    2 & 2 & 1
                \end{bmatrix}, \; B = \begin{bmatrix}
                    2 & 1 \\
                    1 & 3 \\
                    1 & 1
                \end{bmatrix} \; \Longrightarrow \; AB = \begin{bmatrix}
                    5 & 6 \\ 7 & 9
                \end{bmatrix}, \; BA = \begin{bmatrix}
                    4 & 4 & 5 \\
                    7 & 7 & 5 \\
                    3 & 3 & 3
                \end{bmatrix}.
            \end{equation*}
        \end{example}
    \end{frame}

    \begin{frame}{Matrix-Vector Product}
        Vectors are written in columns. For $A \in \mathbb{F}^{m \times n}$ and $\mathbf{x} \in \mathbb{F}^n$:
        \begin{equation*}
            A\mathbf{x} = \begin{bmatrix}
                a_{1,1} & \dots & a_{1,n} \\
                \vdots & \ddots & \vdots \\
                a_{m,1} & \dots & a_{m,n}
            \end{bmatrix}\begin{bmatrix}
                x_1 \\ \vdots \\ x_n
            \end{bmatrix} = \begin{bmatrix}
                a_{1,1}x_1 + \dots + a_{1,n}x_n \\
                \vdots \\
                a_{m,1}x_1 + \dots + a_{m,n}x_n
            \end{bmatrix}.
        \end{equation*}
        \pause

        \begin{remark}
            Why is the matrix product so odd? We study maps $\mathbf{x} \mapsto A\mathbf{x}$. Composing $\mathbf{x} \mapsto B\mathbf{x}$ and then $\mathbf{y} \mapsto A\mathbf{y}$ must again be such a map, $\mathbf{x} \mapsto C\mathbf{x}$, and $C = AB$ is the only way to get it.
        \end{remark}
        \pause

        \begin{questionblock}{Question}
            Let $\mathbf{x} = (1,2,1)$. Compute $\mathbf{x}^{\top}\mathbf{x}$ and $\mathbf{x}\mathbf{x}^{\top}$. For $A = \begin{bmatrix} 1 & 2 \\ 0 & 1 \end{bmatrix}$, compute $A^n$.
        \end{questionblock}
    \end{frame}

    \begin{frame}{Linear Systems and Rank}
        A system of linear equations is just $A\mathbf{x} = \mathbf{t}$:
        \begin{equation*}
            \begin{cases}
                x_1 + 2x_2 + 3x_3 = 2 \\
                2x_1 + x_2 + 3x_3 = 1 \\
                -x_1 + x_2 + x_3 = 0
            \end{cases}
            \Longleftrightarrow \;
            \begin{bmatrix}
                1 & 2 & 3 \\
                2 & 1 & 3 \\
                -1 & 1 & 1
            \end{bmatrix}\begin{bmatrix}
                x_1 \\ x_2 \\ x_3
            \end{bmatrix} = \begin{bmatrix}
                2 \\ 1 \\ 0
            \end{bmatrix}.
        \end{equation*}
        \pause

        \begin{definition}
            The \textbf{rank} of $A = (\mathbf{v}_1,\dots,\mathbf{v}_n)$ is $\text{rank}(A) := \dim \text{span}(\mathbf{v}_1,\dots,\mathbf{v}_n)$. Row rank and column rank are always equal.
        \end{definition}
        \pause

        \begin{proposition}
            For $A \in \mathbb{R}^{n \times n}$ and $\mathbf{t} \in \mathbb{R}^n$:
            \begin{enumerate}
                \item $A\mathbf{x} = \mathbf{t}$ always has a single solution iff $\det A \neq 0$;
                \item \textit{(Rouché-Capelli)} a solution exists iff $\text{rank}(A \,\|\, \mathbf{t}) = \text{rank}(A)$.
            \end{enumerate}
        \end{proposition}
    \end{frame}

    \section{Inner Product}

    \begin{frame}{Inner Product Spaces}
        \begin{definition}
            An \textbf{inner product} on $V$ is a function $\langle \cdot, \cdot \rangle: V \times V \to \mathbb{F}$ that is \textbf{positive} ($\langle \mathbf{v}, \mathbf{v} \rangle \geq 0$), \textbf{definite} ($\langle \mathbf{v}, \mathbf{v} \rangle = 0$ iff $\mathbf{v}=\mathbf{0}$), \textbf{additive}, \textbf{homogeneous} and \textbf{symmetric}.
        \end{definition}
        \pause

        \begin{example}
            The \textbf{dot product} on $\mathbb{F}^n$: $\langle \mathbf{u}, \mathbf{v} \rangle := \mathbf{u}^{\top}\mathbf{v} = \sum_{i=1}^{n} u_i v_i$. For $\mathbf{u}=(1, 2, 3)$, $\mathbf{v} = (2, 4, 3)$, we get $\langle \mathbf{u}, \mathbf{v}\rangle = 2 + 8 + 9 = 19$.
        \end{example}
        \pause

        \begin{definition}
            The \textbf{norm} of $\mathbf{v}$ is $\|\mathbf{v}\| := \sqrt{\langle \mathbf{v}, \mathbf{v} \rangle}$. Vectors $\mathbf{v},\mathbf{u}$ are \textbf{orthogonal} if $\langle \mathbf{v}, \mathbf{u} \rangle = 0$.
        \end{definition}

        \begin{remark}
            Over finite fields or $\mathbb{Z}_N$, ``$\geq 0$'' is meaningless, so we only require additivity, homogeneity and symmetry.
        \end{remark}
    \end{frame}

    \begin{frame}{Inner Products: More Examples}
        \begin{example}
            A symmetric matrix $Q \in \mathbb{R}^{n \times n}$ is \textit{positive-definite} if $\mathbf{x}^{\top}Q\mathbf{x} > 0$ for any non-zero $\mathbf{x}$. Then $\langle \mathbf{x}, \mathbf{y} \rangle_Q := \mathbf{x}^{\top}Q\mathbf{y}$ is an inner product.
        \end{example}
        \pause

        \begin{example}
            Polynomials have ``lengths'' too: on $V := \mathbb{R}[T]^{\leq 2}$, define $\langle p,q \rangle := \int_0^1 p(\tau)q(\tau)\,d\tau$.
        \end{example}
        \pause

        \begin{questionblock}{Question}
            Compute the norm $\|T\|$ in $V$ with the inner product above. Find any $q \in V$ orthogonal to $T^2$.
        \end{questionblock}
    \end{frame}

    \section{Hadamard and Outer Products}

    \begin{frame}{Hadamard Product}
        \begin{definition}
            For $A,B \in \mathbb{F}^{m \times n}$, the \textbf{Hadamard product} $A \odot B$ is the matrix $C$ with $C_{i,j} = A_{i,j}B_{i,j}$: we multiply element-wise.
        \end{definition}
        \pause

        \begin{example}
            \begin{equation*}
                \begin{bmatrix}
                    1 & 1 & 2 \\
                    3 & 0 & 3
                \end{bmatrix} \odot \begin{bmatrix}
                    3 & 2 & 1 \\ 0 & 2 & 1
                \end{bmatrix} = \begin{bmatrix}
                    3 & 2 & 2 \\
                    0 & 0 & 3
                \end{bmatrix}.
            \end{equation*}
        \end{example}
        \pause

        \begin{remarkblock}{Why it matters}
            To say that $\mathbf{z} \in \mathbb{F}^n$ is a bit-vector (every $z_i \in \{0,1\}$), write $\mathbf{z} \odot (\mathbf{1}_n - \mathbf{z}) = \mathbf{0}$. R1CS is also a Hadamard-product relation.
        \end{remarkblock}
    \end{frame}

    \begin{frame}{Outer Product}
        \begin{definition}
            For $\mathbf{u} \in \mathbb{F}^n$, $\mathbf{v} \in \mathbb{F}^m$, the \textbf{outer product} is $\mathbf{u} \otimes \mathbf{v} := \mathbf{u}\mathbf{v}^{\top} \in \mathbb{F}^{n \times m}$, with entries $u_iv_j$.
        \end{definition}
        \pause

        \begin{example}
            \begin{equation*}
                \begin{bmatrix}
                    1 \\ 2 \\ 3
                \end{bmatrix} \otimes \begin{bmatrix}
                    2 \\ 4 \\ 3
                \end{bmatrix} = \begin{bmatrix}
                    1 \\ 2 \\ 3
                \end{bmatrix}\begin{bmatrix}
                    2 & 4 & 3
                \end{bmatrix} = \begin{bmatrix}
                    2 & 4 & 3 \\
                    4 & 8 & 6 \\
                    6 & 12 & 9
                \end{bmatrix}.
            \end{equation*}
        \end{example}
        \pause

        \begin{lemma}
            The outer product is distributive, commutes with scalars, $\mathbf{u} \otimes \mathbf{v} = (\mathbf{v} \otimes \mathbf{u})^{\top}$, and for non-zero $\mathbf{u}, \mathbf{v}$ it is a rank-1 matrix.
        \end{lemma}
    \end{frame}


    \makeclosingframe
\end{document}
